% GENERATED FILE. Edit canonical lessons, then run npm run generate:books.
% canonical-source-hash: fnv1a64:04beeb8f59cf1923
% canonical-lessons: KA-C146-ketta, KA-C146-bhayankara, KA-C146-dhairyavanta, KA-C146-cinnada, KA-C146-belliya

\chapter{Bad, Terrible, Brave, Golden, Silver}
\label{ch:ka-bad-terrible-brave-golden-silver}
\hlchaptermodality{146}

\begin{chapteropening}
\textbf{By the end of this chapter:} \emph{I can say bad, terrible, brave, golden and silver.}

Complete the last lesson of chapter 146: i can say bad, terrible, brave, golden and silver.
\end{chapteropening}

\section[\texorpdfstring{\kn{ಕೆಟ್ಟ} (keṭṭa)}{ಕೆಟ್ಟ (keṭṭa)}]{\kn{ಕೆಟ್ಟ} (keṭṭa) — bad, spoiled}

\label{lesson:KA-C146-ketta}



\begin{quote}
Before the new one: say the Kannada for lazy, then the Kannada for clever.
\end{quote}

\subsection*{You'll want to know: \kn{ಕೆಟ್ಟ}}
\textbf{\kn{ಕೆಟ್ಟ}} — \emph{keṭṭa} — \textquotedblleft{}bad, spoiled\textquotedblright{}.

\begin{grammarlens}[title={What you've built}]
One more word for this chapter.
\end{grammarlens}

\subsection*{Your turn}
\begin{itemize}
\raggedright
  \item \emph{Say it:} \emph{keṭṭa}
  \item \emph{Say it:} \emph{keṭṭa}, once more
  \item \emph{Say it:} say \emph{keṭṭa}
  \item \emph{Recall:} say the Kannada for lazy, then the Kannada for clever, then say \emph{keṭṭa} again
\end{itemize}

\begin{tcolorbox}[breakable,colback=teal!4,colframe=teal!35!black,title={Before you move on}]
What does \kn{ಕೆಟ್ಟ} mean? (\textquotedblleft{}Bad, spoiled\textquotedblright{}.) Say it once more.
\end{tcolorbox}
\section[\texorpdfstring{\kn{ಭಯಂಕರ} (bhayaṅkara)}{ಭಯಂಕರ (bhayaṅkara)}]{\kn{ಭಯಂಕರ} (bhayaṅkara) — terrible}

\label{lesson:KA-C146-bhayankara}



\begin{quote}
Before the new one: say the Kannada for clever, then the Kannada for bad.
\end{quote}

\subsection*{You'll want to know: \kn{ಭಯಂಕರ}}
\textbf{\kn{ಭಯಂಕರ}} — \emph{bhayaṅkara} — \textquotedblleft{}terrible\textquotedblright{}.

\begin{grammarlens}[title={What you've built}]
One more word for this chapter.
\end{grammarlens}

\subsection*{Your turn}
\begin{itemize}
\raggedright
  \item \emph{Say it:} \emph{bhayaṅkara}
  \item \emph{Say it:} \emph{bhayaṅkara}, once more
  \item \emph{Say it:} say \emph{bhayaṅkara}
  \item \emph{Recall:} say the Kannada for clever, then the Kannada for bad, then say \emph{bhayaṅkara} again
\end{itemize}

\begin{tcolorbox}[breakable,colback=teal!4,colframe=teal!35!black,title={Before you move on}]
What does \kn{ಭಯಂಕರ} mean? (\textquotedblleft{}Terrible\textquotedblright{}.) Say it once more.
\end{tcolorbox}
\section[\texorpdfstring{\kn{ಧೈರ್ಯವಂತ} (dhairyavanta)}{ಧೈರ್ಯವಂತ (dhairyavanta)}]{\kn{ಧೈರ್ಯವಂತ} (dhairyavanta) — brave}

\label{lesson:KA-C146-dhairyavanta}



\begin{quote}
Before the new one: say the Kannada for bad, then the Kannada for terrible.
\end{quote}

\subsection*{You'll want to know: \kn{ಧೈರ್ಯವಂತ}}
\textbf{\kn{ಧೈರ್ಯವಂತ}} — \emph{dhairyavanta} — \textquotedblleft{}brave\textquotedblright{}.

\begin{grammarlens}[title={What you've built}]
One more word for this chapter.
\end{grammarlens}

\subsection*{Your turn}
\begin{itemize}
\raggedright
  \item \emph{Say it:} \emph{dhairyavanta}
  \item \emph{Say it:} \emph{dhairyavanta}, once more
  \item \emph{Say it:} say \emph{dhairyavanta}
  \item \emph{Recall:} say the Kannada for bad, then the Kannada for terrible, then say \emph{dhairyavanta} again
\end{itemize}

\begin{tcolorbox}[breakable,colback=teal!4,colframe=teal!35!black,title={Before you move on}]
What does \kn{ಧೈರ್ಯವಂತ} mean? (\textquotedblleft{}Brave\textquotedblright{}.) Say it once more.
\end{tcolorbox}
\section[\texorpdfstring{\kn{ಚಿನ್ನದ} (cinnada)}{ಚಿನ್ನದ (cinnada)}]{\kn{ಚಿನ್ನದ} (cinnada) — golden}

\label{lesson:KA-C146-cinnada}



\begin{quote}
Before the new one: say the Kannada for terrible, then the Kannada for brave.
\end{quote}

\subsection*{You'll want to know: \kn{ಚಿನ್ನದ}}
\textbf{\kn{ಚಿನ್ನದ}} — \emph{cinnada} — \textquotedblleft{}golden\textquotedblright{}.

\begin{grammarlens}[title={What you've built}]
One more word for this chapter.
\end{grammarlens}

\subsection*{Your turn}
\begin{itemize}
\raggedright
  \item \emph{Say it:} \emph{cinnada}
  \item \emph{Say it:} \emph{cinnada}, once more
  \item \emph{Say it:} say \emph{cinnada}
  \item \emph{Recall:} say the Kannada for terrible, then the Kannada for brave, then say \emph{cinnada} again
\end{itemize}

\begin{tcolorbox}[breakable,colback=teal!4,colframe=teal!35!black,title={Before you move on}]
What does \kn{ಚಿನ್ನದ} mean? (\textquotedblleft{}Golden\textquotedblright{}.) Say it once more.
\end{tcolorbox}
\section[\texorpdfstring{\kn{ಬೆಳ್ಳಿಯ} (beḷḷiya)}{ಬೆಳ್ಳಿಯ (beḷḷiya)}]{\kn{ಬೆಳ್ಳಿಯ} (beḷḷiya) — silver}

\label{lesson:KA-C146-belliya}



\begin{quote}
Before the new one: say the Kannada for brave, then the Kannada for golden.
\end{quote}

\subsection*{You'll want to know: \kn{ಬೆಳ್ಳಿಯ}}
\textbf{\kn{ಬೆಳ್ಳಿಯ}} — \emph{beḷḷiya} — \textquotedblleft{}silver\textquotedblright{}.

\begin{grammarlens}[title={What you've built}]
One more word for this chapter.
\end{grammarlens}

\subsection*{Your turn}
\begin{itemize}
\raggedright
  \item \emph{Say it:} \emph{beḷḷiya}
  \item \emph{Say it:} \emph{beḷḷiya}, once more
  \item \emph{Say it:} say \emph{beḷḷiya}
  \item \emph{Recall:} say the Kannada for brave, then the Kannada for golden, then say \emph{beḷḷiya} again
\end{itemize}

\begin{tcolorbox}[breakable,colback=teal!4,colframe=teal!35!black,title={Before you move on}]
What does \kn{ಬೆಳ್ಳಿಯ} mean? (\textquotedblleft{}Silver\textquotedblright{}.) Say it once more.
\end{tcolorbox}
